\chapter{Einleitung}
\label{chap:einleitung}

Bei dem bildbasierten Retrieval ordnet ein System die Bilder einer Bildsammlung danach, wie ähnlich sie einem Anfragebild (Query) sind \cite{musgrave_metric_2020}. Was dabei als ähnlich gilt, hängt jedoch von der Aufgabe ab. Wer ein Foto eines roten Audi A3 von vorn als Query verwendet, sucht weitere Bilder desselben Modells, auch in anderer Farbe und aus anderem Blickwinkel. Ein roter VW Golf von vorn kann der Query stärker ähneln und trotzdem ein falscher Treffer sein.

Bei der Suche nach einem bestimmten Automodell kommt es auf feine Unterschiede zwischen eng verwandten Unterkategorien an. Solche Aufgaben werden als \emph{fine-grained Retrieval} bezeichnet \cite{wei_fine-grained_2021}. Die Suche soll auch für Klassen funktionieren, die im Training nicht vorkamen \cite{song_deep_2016}. Zeigt eine Query etwa einen Audi Q4 e-tron, der nicht zu den Trainingsklassen gehörte, sollen ohne neues Training weitere Bilder dieses Modells gefunden werden. Dazu müssen Merkmale gelernt werden, mit denen sich auch bislang unbekannte Klassen voneinander unterscheiden lassen \cite{song_deep_2016, musgrave_metric_2020}.

Diesen Ansatz verfolgt das Metric Learning. Gelernt wird dabei ein Abstandsmaß, unter dem Bilder derselben Klasse möglichst geringe und Bilder unterschiedlicher Klassen möglichst große Abstände aufweisen sollen \cite{bellet_survey_2014}. Im \ac{DML} berechnet ein tiefes neuronales Netz aus jedem Bild ein Embedding. Beim Retrieval wird dann der Abstand zwischen diesen Embeddings gemessen \cite{musgrave_metric_2020, mohan_deep_2023}. 

Vortrainierte Modelle wie DINOv2 liefern dafür einen guten Ausgangspunkt. DINOv2 wird ohne Labels vortrainiert und schon ein linearer Klassifikator auf seinen eingefrorenen Merkmalen erreicht auf vielen Benchmarks hohe Genauigkeiten \cite{oquab_dinov2_2023}. Solche Merkmale können mit gelabelten Beispielen der Zieldomäne angepasst werden, was die Retrieval-Leistung auf den meisten Datensätzen verbessert \cite{berton_all_2025}. Die Beschaffung solcher Labels ist jedoch aufwendig, da die Unterscheidung eng verwandter Kategorien häufig Fachwissen erfordert \cite{wei_fine-grained_2021}. Daher liegt es nahe, zusätzlich ungelabelte Bilder zu nutzen.


\section{Problemstellung}

\ac{SSL} nutzt oft neben wenigen gelabelten Daten viele ungelabelte Daten \cite{van_engelen_survey_2019}. Wann \ac{SSL}-Verfahren in der Bildklassifikation supervised Baselines übertreffen, ist unter kontrollierten Bedingungen bereits untersucht worden \cite{oliver_realistic_2018, su_realistic_2021, wang_usb_2022, zhang_revisiting_2025}. Die in früheren Arbeiten berichteten Vorsprünge der \ac{SSL}-Verfahren ließen sich bei gleicher Architektur und gleich umfangreicher Hyperparametersuche für alle Verfahren einschließlich der supervised Baseline im Allgemeinen nicht in voller Höhe wiederholen \cite{oliver_realistic_2018}. Benchmarks wie der \ac{USB} vergleichen viele Verfahren unter einheitlichen Bedingungen \cite{wang_usb_2022}.

Auf dem Gebiet des \ac{SSDML} ist die Lage anders. Die in Abschnitt \ref{sec:ssdml} vorgestellten Arbeiten schlagen jeweils ein eigenes Verfahren vor, etwa Triplets aus über einen Graphen propagierten Affinitäten \cite{dutta_semi-supervised_2021}, Label Propagation mit Proxys \cite{zhuang_deep_2021, zhuang_deep_2023}, eine Art von Self-Training \cite{duan_slade_2021, sahito_semi-supervised_2021} oder ein für das Lernen ohne Labels entwickeltes Teacher-Student-Verfahren mit weichen Pseudo-Labels \cite{kim_self-taught_2022}. \textcite{dutta_semi-supervised_2021} verglichen ihr Verfahren direkt mit fünf klassischen semi-supervised Metric-Learning-Verfahren, die sie auf tiefe Netze übertrugen. Dort wurde je Datensatz nur ein Label-Budget untersucht. Die Ergebnisse der supervised Baselines werden nur für \ac{CUB} berichtet. Für alle Versuche wird eine Validierungsmenge zur Abstimmung der Hyperparameter genannt. Die gewählten Hyperparameterwerte werden jedoch nur für das eigene Verfahren angegeben \cite{dutta_semi-supervised_2021}. Die übrigen Arbeiten verglichen das eigene Verfahren mit supervised Baselines, untersuchten aber den Einfluss unterschiedlicher Label-Budgets nicht und beschreiben keine eigene Hyperparametersuche für die Vergleichsverfahren \cite{zhuang_deep_2021, zhuang_deep_2023, duan_slade_2021, kim_self-taught_2022, sahito_semi-supervised_2021}. Teils wurden fremde Ergebnisse übernommen, die unter anderen Einstellungen als die eigenen Versuche erzielt wurden \cite{duan_slade_2021, zhuang_deep_2023}. Die sechs Arbeiten unterscheiden sich außerdem untereinander teilweise im Backbone, in den Datensätzen, in den Label-Budgets und in der Herkunft der ungelabelten Daten \cite{dutta_semi-supervised_2021, zhuang_deep_2021, zhuang_deep_2023, duan_slade_2021, kim_self-taught_2022, sahito_semi-supervised_2021}. Für das supervised \ac{DML} wurde bereits gezeigt, dass die Unterschiede zwischen Verfahren unter einheitlichen Bedingungen schrumpfen \cite{musgrave_metric_2020}. Die genannten Arbeiten zeigen zwar Verbesserungen durch ungelabelte Bilder, lassen aber offen, welche Mechanismen unter welchen Bedingungen einen zusätzlichen Nutzen bringen.

Diese Frage stellt sich besonders bei stark vortrainierten Modellen. Auf dem fine-grained Vogeldatensatz Semi-Aves hebt ein Vortraining auf ImageNet die Top-1-Genauigkeit der supervised Baseline von \qty{20.6}{\percent} auf \qty{52.7}{\percent}, das beste \ac{SSL}-Verfahren auf demselben vortrainierten Netz erreicht \qty{57.4}{\percent} \cite{su_realistic_2021}. Auch bei vortrainierten \ac{ViT}-Modellen ist der Abstand zwischen \ac{SSL}-Verfahren und einer supervised Baseline mit parametereffizientem Fine-Tuning oft gering oder kehrt sich sogar um \cite{zhang_revisiting_2025}. Für diese Arbeit folgt daraus die Vermutung, dass ein großer Teil des möglichen Beitrags ungelabelter Bilder bereits im Vortraining steckt. Hinzu kommt, dass semi-supervised Verfahren auf Annahmen über die Struktur der Daten beruhen, etwa dass nahe beieinanderliegende Datenpunkte derselben Klasse angehören \cite{van_engelen_survey_2019}. Für das \ac{DML} bedeutet das, dass solche Annahmen für die Trainingsbilder gelten müssen, während auf Klassen bewertet wird, die im Training fehlen. Ob sich die Befunde aus der Klassifikation auf diese Bedingungen übertragen lassen, ist daher offen.

Die Frage wird anhand der Datensätze Cars196 mit Automodellen, \ac{CUB} mit Vogelarten und In-Shop mit Kleidungsstücken untersucht. Alle drei Datensätze werden häufig zum Vergleich von \ac{DML}-Verfahren verwendet \cite{song_deep_2016, musgrave_metric_2020, wang_multi-similarity_2019, kim_proxy_2020}. Der vierte Datensatz, Semi-iNat-2021, wurde für die Evaluation von semi-supervised Verfahren zusammengestellt und enthält viele Arten, deren Bildzahlen stark ungleich verteilt sind \cite{su_semi-supervised_inat_2021}. Mit ihm wird die Frage unter anspruchsvolleren Bedingungen geprüft. Der Beitrag ungelabelter Bilder könnte auch davon abhängen, wie gut die vortrainierten Merkmale einen Datensatz bereits erfassen. In die Vortrainingsdaten von DINOv2 wurden gezielt Bilder aufgenommen, die den Trainingsbildern von Cars196 und \ac{CUB} ähneln \cite{oquab_dinov2_2023}.

\section{Zielsetzung und Forschungsfragen}

Diese Arbeit vergleicht vier Mechanismen, mit denen ungelabelte Bilder in das Training eines Modells eingehen. Sie stehen für das graphbasierte Pseudo-Labeling, die Mannigfaltigkeitsregularisierung, die Entropieminimierung und das clusterbasierte Pseudo-Labeling als vier Richtungen des \ac{SSL}.

\emph{\ac{LP}} überträgt Labels über einen Graphen auf ungelabelte Bilder \cite{zhu_learning_2002, zhou_learning_2003}. Beim Training tiefer Netze dienen die übertragenen Labels als Pseudo-Labels, mit denen das Netz trainiert wird, bevor aus seinen Embeddings der nächste Graph entsteht \cite{iscen_label_2019}. \emph{\ac{LRML}} bestraft über einen Graph-Laplace-Term große Abstände zwischen Bildern, die in einem \aclu{kNN}-Graphen (\acs{kNN}-Graphen) benachbart sind \cite{hoi_semi-supervised_2010}. Die \emph{Entropieminimierung nach Hoffer und Ailon} fördert, dass jedes ungelabelte Bild einem der gelabelten Referenzbilder deutlich näher liegt als den übrigen \cite{hoffer_semi-supervised_2018}. \emph{\ac{SLADE}} leitet aus Clustern eines Teachers Pseudo-Labels ab \cite{duan_slade_2021}.

Jeder Mechanismus wurde mit fünf supervised Verlustfunktionen kombiniert und mit einer supervised Baseline derselben Verlustfunktion verglichen. Trainiert wurde nur ein linearer Projection Head auf seinen einmal berechneten Merkmalen. Dadurch war jeder Trainingslauf kurz genug, um alle Kombinationen aus Mechanismus, Verlustfunktion und Datensatz bei einem Label-Budget mit eigener Hyperparametersuche zu untersuchen.

Die Arbeit wird von drei Forschungsfragen geleitet. Die erste Frage (F1) lautet, ob ungelabelte Bilder die Retrieval-Leistung auf ungesehenen Testklassen gegenüber einer supervised Baseline mit gleich umfangreicher Hyperparametersuche verbessern, wenn die Merkmale von einem vortrainierten \ac{ViT} stammen. Die zweite Frage (F2) lautet, wie ein solcher Nutzen vom Label-Budget, vom Datensatz und von der Verlustfunktion abhängt und worauf er beruht. Die dritte Frage (F3) lautet, welche Hyperparameter und Designentscheidungen bestimmen, ob ein Mechanismus hilft.

Zu den Forschungsfragen F1 und F2 werden drei Hypothesen geprüft, die aus den Befunden zur semi-supervised Klassifikation abgeleitet wurden. H1 besagt, dass ungelabelte Bilder die Retrieval-Leistung gegenüber der supervised Baseline verbessern. H2 besagt, dass dieser Nutzen bei kleinen Label-Budgets größer ist als bei großen. H1 und H2 beruhen darauf, dass ungelabelte Daten zusätzliche Informationen liefern können, vor allem wenn nur wenigen Beispielen ein Label zugewiesen ist \cite{van_engelen_survey_2019}. Sie geben damit die übliche Erwartung an das \ac{SSL} wieder, die die Vermutung aus der Problemstellung für stark vortrainierte Modelle infrage stellt. H3 besagt, dass der Nutzen sinkt, wenn die ungelabelten Bilder aus Klassen ohne gelabelte Beispiele stammen. In der Klassifikation verschlechterten ungelabelte Fahrzeugbilder neben gelabelten Tierklassen die Ergebnisse auf CIFAR-10 \cite{oliver_realistic_2018}. Wurden die Bilder fremder Klassen nur hinzugefügt, ohne ungelabelte Bilder bekannter Klassen zu entfernen, blieb dieser Schaden teils aus \cite{he_reevaluating_2025}. Bei vortrainierten Modellen schadeten dagegen häufig auch Bilder weiterer Vogel- und Pilzarten, die zu den ungelabelten Bildern der gelabelten Arten hinzukamen \cite{su_realistic_2021}. Erwartet wird dieser Rückgang bei \ac{LP} und der Entropieminimierung, weil nur sie jedes ungelabelte Bild einer gelabelten Klasse zuordnen.

Den zweiten Teil von F2, worauf ein Nutzen beruht, prüfen Ablationen. Ein Nutzen könnte von den ungelabelten Bildern selbst stammen, von den eigens gesuchten Hyperparametern der \ac{SSL}-Konfiguration oder vom \ac{SSL}-Mechanismus selbst. Die erwartete Wirkung der Ablationen wurde mit einer Ausnahme vor den Experimenten festgelegt. Zusammen mit Sweeps, die jeweils einen Hyperparameter variieren, bilden die Ablationen auch die wichtigste Grundlage für die Antwort auf F3.

\section{Beiträge und Aufbau der Arbeit}

Erstens liefert die Arbeit einen Vergleich der vier genannten Mechanismen mit je fünf Verlustfunktionen auf den Datensätzen Cars196, \ac{CUB}, In-Shop und Semi-iNat-2021 bei mehreren Label-Budgets. Jede Kombination aus Datensatz, Mechanismus und Verlustfunktion erhielt, ebenso wie die supervised Baseline, eine eigene Hyperparametersuche über 80 Versuche, davon 40 mit dem \ac{TPE}, bewertet mit klassendisjunkter Kreuzvalidierung. Zweitens führen Ablationen einen Nutzen der Mechanismen auf seine Ursachen zurück und prüfen, welche Rolle die Qualität der Pseudo-Labels oder des Graphen spielt. Drittens wurde die Geometrie des Embedding-Raums über den Trainingsverlauf verfolgt. Geprüft wurde, ob sich der Zusammenhang zwischen Geometriemaßen und der Leistung auf ungesehenen Klassen \cite{roth_revisiting_2020,milbich_characterizing_2021} auch in diesem Versuchsaufbau zeigt. Viertens wurde untersucht, wie der Nutzen ungelabelter Bilder von der Qualität der vortrainierten Merkmale auf dem jeweiligen Datensatz abhängt.

In Kapitel \ref{chap:grundlagen} werden die Grundlagen des \ac{DML} eingeführt. Dazu gehören die Grundbegriffe des Metric Learning, die Verlustfunktionen und das Mining sowie die Bewertung anhand von Retrieval-Metriken und Maßen für die Geometrie eines Embedding-Raums. Außerdem werden das Backbone DINOv2 und die Hyperparameteroptimierung vorgestellt. In Kapitel \ref{chap:ssdml} werden die Annahmen, die Evaluation und die \ac{SSL}-Verfahren sowie deren Übertragung auf das Metric Learning behandelt. Dazu gehören auch die in dieser Arbeit verwendeten Mechanismen. In Kapitel \ref{chap:methodik} werden der Versuchsaufbau, die Datensätze, das Auswertungsprotokoll und die Umsetzung der Mechanismen beschrieben. Das Kapitel erläutert außerdem das Vorgehen von den Vorexperimenten über die Hyperparameteroptimierung (Phase 1) bis zu den Ablationen und Analysen (Phase 2). Schließlich werden in Kapitel \ref{chap:ergebnisse} die Ergebnisse der Experimente präsentiert und in Kapitel \ref{chap:diskussion} eingeordnet.